\documentclass[11pt]{article}
% Préambule commun des prétirages v2 (format arXiv : source TeX autonome, pdfLaTeX).
% Moteur : pdfLaTeX (arXiv) → fontenc T1 + inputenc ; XeTeX/LuaTeX (tectonic) → fontspec. Sans ce choix, les
% guillemets « » tapés en dur sortent en « ń ż » sous tectonic (ledger, 2026-11-14).
\usepackage{iftex}
\ifPDFTeX
  \usepackage[T1]{fontenc}
  \usepackage[utf8]{inputenc}
\else
  \usepackage{fontspec}
\fi
\usepackage[french]{babel}
\usepackage{lmodern,microtype}
\usepackage[a4paper,margin=2.5cm]{geometry}
\usepackage{amsmath,amssymb}
\usepackage{graphicx,booktabs,siunitx}
\usepackage{authblk}
\usepackage[numbers,sort&compress]{natbib}
\usepackage[hidelinks]{hyperref}
\usepackage{csquotes}
\sisetup{locale=FR,group-minimum-digits=4}
\renewcommand\Affilfont{\small}
% Étiquette de provenance obligatoire sur chaque nombre (règle du projet) :
\newcommand{\mesure}{\textsuperscript{\textsc{m}}}   % mesuré
\newcommand{\derive}{\textsuperscript{\textsc{d}}}   % dérivé
\newcommand{\suppose}{\textsuperscript{\textsc{s}}}  % supposé


\title{Vingt-sept hyperparamètres sourcés pour un jumeau numérique de l'offre de soins \enquote{libérale}~: catalogue, statut de preuve et lien aux leviers de politique publique}
\author[1]{Équipe hackathon-santé}
\affil[1]{Projet \enquote{jumeau numérique de l'offre de soins}, dépôt \texttt{hackathon-sante/v2}}
\date{Prétirage, version du \today{} --- non relu par des pairs}

\begin{document}
\maketitle

\begin{abstract}
Un jumeau numérique de l'offre de soins simule des politiques publiques en
réglant des hyperparamètres. Encore faut-il que chacun ait une valeur, une
source primaire, un millésime et une incertitude. Nous dressons ce catalogue
pour cinq professions (sages-femmes, chirurgiens-dentistes, pharmaciens,
masseurs-kinésithérapeutes, infirmiers) sur quatre mécanismes propres à chaque
profession : régulation conventionnelle de l'installation, filières en
pharmacie, spécialisation en odontologie, abandon en études infirmières.
Nous y ajoutons les paramètres de sortie communs.

Le catalogue compte 27~hyperparamètres : 19~sourcés, 1~mesuré, 1~supposé et
6~non mesurables à ce jour. Un seul effet causal de la régulation est
publié : $\num{-25.57}$ infirmières libérales pour \num{100000}~habitants en
zone sur-dotée, par rapport aux zones intermédiaires~\cite{irdes270}.

Deux affirmations courantes ne résistent pas aux sources.
\begin{itemize}
\item Le \enquote{30~\% d'abandon} en institut de formation en soins
  infirmiers (IFSI) n'a pas de source primaire. Ce que la DREES mesure, c'est
  un rendement diplômés / entrants passé de \num{0.84}\derive{} à
  \num{0.72}\derive{}.
\item L'obligation d'exercer deux ans en zone sous-dotée vise les
  kinésithérapeutes, et non les dentistes.
\end{itemize}
Nous traduisons enfin 14~politiques publiques en réglages de ces
hyperparamètres, pour l'interface du décideur. Huit ont un point d'entrée
dans le code du jumeau ; trois sont jouables aujourd'hui.
Limite principale : aucun effet chiffré de la régulation n'existe pour les
sages-femmes, les kinésithérapeutes ou les dentistes, ni sur la destination
des entrants bloqués.
\end{abstract}
% Le préambule commun ne charge que babel[french] : l'environnement otherlanguage{english}
% du gabarit échouerait à la compilation. On renomme seulement le titre du résumé.
{\renewcommand{\abstractname}{Abstract}
\begin{abstract}
A digital twin of the French health workforce simulates public policies by
setting hyperparameters. Each one needs a value, a primary source, a vintage
and an uncertainty. We build this catalogue for five professions (midwives,
dentists, pharmacists, physiotherapists, nurses). It covers four
profession-specific mechanisms: contractual regulation of new practices,
career tracks in pharmacy, dental specialisation, and attrition in nursing
school. It also covers the exit parameters shared by all professions.

The catalogue holds 27 hyperparameters: 19 sourced, 1 measured, 1 assumed
and 6 not yet measurable. Only one causal estimate of regulation is
published: $-25.57$ self-employed nurses per 100{,}000 inhabitants in
over-supplied areas relative to intermediate areas~\cite{irdes270}.

Two common claims fail against the sources. The ``30\,\% attrition'' figure
in nursing school has no primary source; what the DREES data measure is a
graduates-to-entrants ratio that fell from 0.84 to 0.72. The two-year
obligation to practise in under-served areas applies to physiotherapists,
not to dentists.

Finally, we map 14 public policies onto these hyperparameters for the
decision-maker's interface. Eight have an entry point in the twin's code;
three can be played today.

Main limitation: no quantified effect of regulation exists for midwives,
physiotherapists or dentists, nor on where blocked entrants go.
\end{abstract}}
\noindent\textbf{Mots-clés :} démographie des professions de santé ;
jumeau numérique ; régulation de l'installation ; zonage conventionnel ;
abandon en formation ; hyperparamètres sourcés.

\paragraph{Convention de provenance.}
Le préambule commun définit trois étiquettes :
\begin{itemize}
\item \mesure{} : calculé par nous sur une donnée ouverte téléchargée ;
\item \derive{} : rapport ou différence de nombres publiés, calculé par nous ;
\item \suppose{} : valeur par défaut sans mesure.
\end{itemize}
Un nombre suivi d'une référence entre crochets est \emph{cité} : il
correspond au statut \enquote{sourcé} des fiches. Il est lu dans la source
primaire indiquée.

\section{Introduction}

\paragraph{Le projet.}
Le projet modélise l'offre de soins \enquote{libérale} à l'horizon 2050, de
la commune au niveau national.
\begin{itemize}
\item \emph{Libérale} s'entend au sens du secteur privé, lucratif ou non.
\item Tout employeur public est exclu.
\end{itemize}
Une politique publique y est un scénario, c'est-à-dire un réglage
d'hyperparamètres. Les modèles de référence ne suffisent pas.
\begin{itemize}
\item Omphale (Insee) projette la population.
\item Les projections de la DREES par profession~\cite{dd76,er1062,er1075,er1319}
  ne paramètrent que la formation, la cessation et les diplômés étrangers.
\end{itemize}
Aucun des deux n'expose les mécanismes propres à chaque profession sur
lesquels portent les politiques récentes : régulation conventionnelle de
l'installation, internat, abandon en formation.

\paragraph{Question.}
Quels hyperparamètres propres à chaque profession peut-on adosser à une
source primaire chiffrée et datée ? Lesquels restent, à ce jour, non
mesurables ?

\paragraph{Contributions.}
\begin{enumerate}
\item Un catalogue de 27~fiches. Chacune donne un mécanisme, une valeur,
  une unité, une incertitude, une source primaire, un millésime et un statut
  de preuve. Les fiches sont dans \texttt{v2/docs/lab/fiches/param-*.md}.
\item Pour la régulation, les textes conventionnels exacts des quatre
  professions concernées (règle, zones, dates d'effet). Ils sont recoupés
  avec le zonage ouvert par commune.
\item Deux corrections de prémisses :
  \begin{itemize}
  \item le chiffre de 30~\% d'abandon en IFSI n'a pas de source primaire ;
  \item l'obligation temporaire d'exercice n'existe pas pour les dentistes.
  \end{itemize}
\item Un rendement de formation infirmière reproductible à partir de la
  série longue de la DREES, avec son script.
\item La correspondance entre chaque hyperparamètre, le module du jumeau qui
  le lit et le levier de politique publique qui le déplace.
\end{enumerate}

\section{Travaux existants}

\paragraph{Projections de la DREES.}
Le Dossier n\textsuperscript{o}~76~\cite{dd76} projette les dentistes, les
sages-femmes et les pharmaciens jusqu'en 2050 avec un modèle par
composantes.
\begin{itemize}
\item Les groupes sont définis par âge, sexe et mode d'exercice.
\item La cessation y est la probabilité d'être actif une année et inactif la
  suivante. Elle est estimée sur le panel du RPPS 2016--2020.
\item Les valeurs de cette probabilité ne sont pas publiées.
\end{itemize}
Les infirmiers et les kinésithérapeutes ont été projetés en
2018~\cite{er1062,er1075}, puis les infirmiers mis à jour en 2024
jusqu'en 2050~\cite{er1319}.
\begin{itemize}
\item Ces travaux bornent l'exercice à 67~ans (infirmiers) et à 70~ans
  (kinésithérapeutes), et n'ont aucune cessation avant 40~ans~\cite{er1062,er1075}.
\item Aucun ne sépare la mortalité de la cessation.
\end{itemize}

\paragraph{Évaluation de la régulation.}
L'IRDES~\cite{irdes270} estime par différence de différences l'effet du
zonage infirmier de 2012 sur la période 2006--2016. C'est la seule
estimation causale chiffrée que nous ayons trouvée.
\begin{itemize}
\item L'IGAS~\cite{igas2023} constate une baisse des disparités pour les
  infirmiers et les sages-femmes entre 2016 et 2019, sans en donner
  l'ampleur.
\item La Cour des comptes~\cite{ccomptes2024} relève un report des
  installations en limite des zones sur-denses.
\item Un mémoire de l'EHESP~\cite{bahroun2025} recense les dispositifs de
  régulation par profession et cartographie les zones concernées.
\end{itemize}

\paragraph{Abandon en IFSI.}
L'étude de la DREES de 2023~\cite{er1266} mesure 10~\% d'abandon en
première année en 2021, contre 3~\% en 2011. Elle mesure aussi 14~\%
d'abandon sur toute la scolarité pour la promotion entrée en 2018.
L'étude de 2024~\cite{er1319} montre l'équivalence de deux leviers pour
maintenir l'offre infirmière :
\begin{itemize}
\item diviser par deux le \enquote{taux de perte} en formation ;
\item augmenter durablement de 14~\% le nombre de places.
\end{itemize}

\section{Données}

Le tableau~\ref{tab:donnees} recense les jeux de données utilisés. Tous
sont ouverts ou publiés ; le millésime retenu est indiqué.

\begin{table}[htbp]
\centering\small
\caption{Sources primaires du catalogue.}\label{tab:donnees}
\begin{tabular}{p{4.2cm}p{3.2cm}p{3cm}p{3.6cm}}
\toprule
Source & Producteur & Millésime & Grain \\
\midrule
Avenants et convention~\cite{avenant6ide,avenant4sf,avenant7mk,conventioncd2023} & UNCAM et syndicats & 2018--2023 & national (règle) \\
Zonage par commune~\cite{zonagesdatagouv} & Atlasanté/DataSanté & 20/10/2026 & commune, COG 2026 \\
Panorama démographique~\cite{cnop2025,cnop2024,cnop2016} & Ordre des pharmaciens & 1/1/2026, 1/1/2025, 2015 & national, section \\
Postes d'internat~\cite{cng2018,arrete2023odonto,arrete2024odonto} & CNG, JORF & 2008--2024 & national, DES \\
Rapports d'activité~\cite{oncd2022} & Ordre des chirurgiens-dentistes & 2019--2022 & national, spécialité \\
Enquête Écoles~\cite{drees491} & DREES & 1966--2025 & national (métropole et DROM) \\
Démographie des kinésithérapeutes~\cite{cnomk2024} & CNOMK & 2019--2023 & national \\
Cumul emploi-retraite~\cite{cor2016} & COR (données des caisses) & 30/06/2015 & caisse \\
\bottomrule
\end{tabular}
\end{table}

\paragraph{Contrôles.}
\begin{itemize}
\item \textbf{Zonage.} Le fichier data.gouv concorde avec l'extraction
  CartoSanté de l'axe~4 sur le nombre de communes en \enquote{remplacement
  seul}, à dix communes près\derive. Exemple pour les kinésithérapeutes :
  8,1~\% de \num{35002}~communes, soit \num{2835}\derive{}, contre
  \num{2836}\mesure{} géocodes.
\item \textbf{Graphiques de l'Ordre des pharmaciens.} L'extraction texte
  des camemberts inverse les étiquettes des sections~A et~B. Les parts ont
  donc été relues sur le rendu graphique.
\end{itemize}

\section{Méthode}

\paragraph{Fiche.}
Chaque hyperparamètre a une fiche qui donne :
\begin{itemize}
\item le module du jumeau qu'il alimente : formation, abandon, filières,
  installation, régulation, sorties ou transitions ;
\item sa valeur, son unité et son incertitude ;
\item sa source primaire et son millésime ;
\item un statut parmi quatre : \emph{sourcé} (lu dans une source
  primaire), \emph{mesuré} (calculé par nous sur une donnée ouverte),
  \emph{supposé} ou \emph{non mesurable}.
\end{itemize}
Une piste sans source chiffrée est classée \enquote{non mesurable}, jamais
complétée par une valeur inventée. La fiche dit alors comment la mesurer et
à quel coût.

\paragraph{Recherche.}
Six recherches documentaires indépendantes ont couvert : les textes de la
régulation, ses effets, les filières en pharmacie, l'odontologie, l'IFSI et
les sorties. Une source n'entre au catalogue que si son texte a été ouvert
et lu. Une reprise par la presse sert au mieux de confirmation.

\paragraph{Formes fonctionnelles.}
Ces formes disent comment le jumeau lit chaque paramètre ; elles ne sont pas
estimées ici. Notations :
\begin{itemize}
\item $p$ : la profession ; $c$ : la commune ; $t$ : l'année ;
\item $E$ : les entrées en exercice libéral ; $X$ : les cessations.
\end{itemize}

\emph{Régulation.} Dans une zone régulée, le conventionnement exige la
succession d'un départ :
\begin{equation}
E_p(c,t) \le X_p(c,t), \qquad c \in \mathcal{Z}_p^{\text{régulée}}.
\end{equation}
Un plafond qui fuit, de paramètre $\theta_p \in [0,1]$, s'écrit :
\begin{equation}
E_p(c,t) = \theta_p \min\bigl(E^{\text{libre}}_p, X_p\bigr) + (1-\theta_p)\,E^{\text{libre}}_p .
\end{equation}
Les entrants bloqués, $E^{\text{libre}}_p - E_p$, sont ventilés par un
vecteur $\psi_{p,d}$ entre cinq destinations : commune voisine,
remplacement, salariat privé, salariat public, sortie.

\emph{Filières en pharmacie.} Les primo-inscrits de l'année se répartissent
entre les sections de l'Ordre selon un vecteur $\varphi_{\mathrm{PH},f}$,
de somme~1.
\begin{itemize}
\item Les sections~A, B, C et~D sont privées.
\item Les sections~G et~H sont partagées entre privé et public, selon des
  parts privées $s_G$ et $s_H$.
\end{itemize}

\emph{Formation infirmière.} Le rendement est
\begin{equation}
\rho_{\mathrm{IDE}}(t) = \frac{D(t+3)}{I_1(t)},
\end{equation}
où $D$ désigne les diplômés et $I_1$ les inscrits en première année. Il
englobe abandon, redoublement et passerelles ; $1-\rho$ n'est donc pas un
taux d'abandon.

\emph{Sorties.} La sortie par âge $x_p(a)$ suit la définition de la
DREES~\cite{dd76}. Elle inclut décès, retraite et départs.

\section{Résultats}

\subsection{Catalogue}

Le tableau~\ref{tab:catalogue} donne les 27~hyperparamètres. Répartition
par statut : 19~sourcés, 1~mesuré, 1~supposé, 6~non mesurables.

\begin{table}[p]
\centering\scriptsize
\caption{Catalogue des hyperparamètres (fiches \texttt{param-*}). Prof. : SF sages-femmes, CD dentistes, PH pharmaciens, MK kinésithérapeutes, IDE infirmiers. Statut : So sourcé, Me mesuré, Su supposé, NM non mesurable.}\label{tab:catalogue}
\begin{tabular}{p{2.6cm}p{1.1cm}p{1.6cm}p{6.6cm}p{0.5cm}}
\toprule
Fiche & Prof. & Module & Valeur (millésime) & Statut \\
\midrule
\multicolumn{5}{l}{\emph{Régulation de l'installation}} \\
ide-regul-regle & IDE & régulation & sur-dotée : 1 entrée pour 1 cessation ; deux tiers de l'activité dans la zone en périphérie ; 27,4~\%\derive{} de la population (avenant 6, 2019~\cite{avenant6ide}) & So \\
ide-regul-effet-densite & IDE & régulation & $\num{-25.57}$ IDEL / \num{100000} hab. ; $\num{-18.8}$ pour les moins de 40~ans (2006--2016~\cite{irdes270}) & So \\
regul-theta-fuite & IDE & régulation & $\theta = \num{0.80}$ $[\num{0.74} ; \num{0.87}]$\suppose{} (axe~3, exploratoire) & Su \\
sf-regul-regle & SF & régulation & sur-dotée : 1 pour 1 ; 12,2~\%\derive{} de la population (avenant 4, 2018~\cite{avenant4sf}) & So \\
sf-regul-effet & SF & régulation & direction seulement (disparités en baisse, 2016--2019~\cite{igas2023}) & NM \\
mk-regul-regle & MK & régulation & non prioritaire : 1 pour 1 ; 26,5~\%\derive{} de la population (avenant 7, 2023~\cite{avenant7mk}) & So \\
mk-obligation-exercice & MK & installation & 2 ans en zone (très) sous-dotée ou 2 ans en établissement ; premiers diplômés concernés en 2027~\cite{avenant7mk} & So \\
mk-regul-effet & MK & régulation & non chiffré & NM \\
cd-regul-regle & CD & régulation & non prioritaire : 1 pour 1 ; 4,8~\%\derive{} de la population ; spécialistes exemptés ; effet au 1/1/2025~\cite{conventioncd2023} & So \\
cd-regul-effet & CD & régulation & non chiffré ; borne haute d'environ 5~\%\derive{} du flux & NM \\
regul-report-entrants & 4 prof. & régulation & $\psi$ : qualitatif seulement (report en bordure, remplacement~\cite{irdes270,ccomptes2024}) & NM \\
\midrule
\multicolumn{5}{l}{\emph{Filières en pharmacie}} \\
ph-filiere-primo-inscrits & PH & filières & D 71,0 ; H 9,5 ; B 8,6 ; G 5,7 ; E 3,5 ; A 1,2 ; C 0,4~\% de \num{3044} primo-inscrits (1/1/2026~\cite{cnop2025}) & So \\
ph-part-privee-biologie & PH & transitions & $s_G$~=~55,8~\%\derive{} (1/1/2026~\cite{cnop2025}) & So \\
ph-part-privee-etablissements & PH & transitions & $s_H$~=~28,6~\%\derive{} des inscriptions (Ordre~\cite{cnop2025}) ; 17~\% selon la DREES (axe~4) & So \\
ph-internat-postes & PH & filières & \num{746} postes (biologie \num{266}, pharmacie hospitalière \num{480}, rentrée 2025~\cite{cnop2024}) ; environ 14~\%\derive{} d'une promotion & So \\
ph-taux-inscription-diplomes & PH & abandon & 32,8~\% de diplômés non inscrits à 3~ans (2015~\cite{cnop2016}) & So \\
\midrule
\multicolumn{5}{l}{\emph{Spécialisation en odontologie}} \\
cd-specialisation-des & CD & filières & 9,2~\%\derive{} d'une promotion ; ODF 52, CO 25, MBD 44 postes (2024--2025~\cite{arrete2024odonto}) & So \\
cd-stock-specialistes & CD & filières & ODF \num{2465}, CO \num{222}, MBD \num{100}, omnipraticiens \num{43242} (2022~\cite{oncd2022}) & So \\
\midrule
\multicolumn{5}{l}{\emph{Abandon en études infirmières}} \\
ide-rendement-diplomes & IDE & formation & $\rho = \num{0.72}$\mesure{} (entrée 2022~\cite{drees491}) & Me \\
ide-abandon-promotion & IDE & abandon & 25,0~\% toutes interruptions (promotion 2025) ; 11,9~\% définitives (promotion 2024)~\cite{drees491} & So \\
ide-abandon-1a & IDE & abandon & 12,0~\% d'interruptions en première année (2025~\cite{drees491}) & So \\
\midrule
\multicolumn{5}{l}{\emph{Sorties et entrée en exercice (toutes professions)}} \\
cessation-par-age & toutes & sorties & non publiée ; mesurable sur le RPPS & NM \\
mortalite & toutes & sorties & incluse dans la cessation~\cite{dd76} & NM \\
age-entree-exercice & toutes & installation & SF 24, CD 26, PH 27 à 29~ans (RPPS 2020~\cite{dd76}) & So \\
cumul-emploi-retraite & toutes & sorties & CD 4,7 ; PH 1,8 ; MK et IDE 1,5 ; SF 0,8~\%\derive{} (2015~\cite{cor2016}) & So \\
mk-age-retraite & MK & sorties & 66,8 à 67,3~ans (2019--2023~\cite{cnomk2024}) & So \\
ph-sorties-ordre & PH & sorties & \num{3902} radiations en 2024 ; 18,4~\% de 60~ans et plus~\cite{cnop2024} & So \\
\bottomrule
\end{tabular}
\end{table}

\subsection{Régulation conventionnelle de l'installation}

\paragraph{Une même règle, des périmètres très différents.}
Les quatre professions appliquent la même règle : en zone sur-dotée ou
\enquote{non prioritaire}, un conventionnement n'est accordé qu'en
succession d'une cessation~\cite{avenant6ide,avenant4sf,avenant7mk,conventioncd2023}.
La part de la population couverte varie d'un facteur six\derive :
\begin{itemize}
\item dentistes : 4,8~\%\derive ;
\item sages-femmes : 12,2~\%\derive ;
\item kinésithérapeutes : 26,5~\%\derive ;
\item infirmiers : 27,4~\%\derive.
\end{itemize}
Ces parts sont calculées sur \num{35002}~communes et 68,9~millions
d'habitants~\cite{zonagesdatagouv}.

\paragraph{Spécificités par profession.}
\begin{itemize}
\item \textbf{Infirmiers.} Un nouvel installé en périphérie d'une zone
  sur-dotée doit réaliser les deux tiers de son activité dans sa zone
  d'installation~\cite{avenant6ide}.
\item \textbf{Dentistes.} Les spécialistes (orthopédie dento-faciale,
  chirurgie orale, médecine bucco-dentaire) sont exemptés~\cite{conventioncd2023}.
\item \textbf{Kinésithérapeutes.} Ceux qui entrent en formation à partir
  de 2023 doivent, pour être conventionnés, choisir entre deux options~\cite{avenant7mk} :
  \begin{itemize}
  \item exercer leurs deux premières années en zone sous-dotée ou très
    sous-dotée ;
  \item justifier de deux ans d'exercice en établissement.
  \end{itemize}
  Le Conseil d'État a rejeté les recours contre cette règle~\cite{ce488665}.
  La note de cadrage du projet l'attribuait à tort aux dentistes.
\end{itemize}

\paragraph{Effet mesuré.}
L'IRDES mesure, après 2012~\cite{irdes270} :
\begin{itemize}
\item $\num{-25.57}$ infirmières libérales pour \num{100000}~habitants en
  zone sur-dotée par rapport aux zones intermédiaires ;
\item $\num{-18.8}$ chez les moins de 40~ans ;
\item $+\num{4.31}$ en zone très sous-dotée, effet attribué aux
  incitations.
\end{itemize}
Les zones témoins reçoivent vraisemblablement une partie du report. L'effet
net national est donc plus faible que ce contraste.

L'axe~3 du projet mesure $-30~\%$\mesure{} $[-34 ; -26]$ de densité
relative sur 2012--2019. Le test de tendances parallèles pré-enregistré
échoue cependant ($p = \num{0.043}$\mesure{}). La valeur $\theta = \num{0.80}$
qui en découle reste donc \emph{supposée}\suppose{}.

\paragraph{Ce qui manque.}
\begin{itemize}
\item Aucun effet n'est chiffré pour les sages-femmes, les kinésithérapeutes
  ou les dentistes~\cite{igas2023,ccomptes2024}.
\item Aucune source ne chiffre la destination des entrants bloqués $\psi$.
\end{itemize}

\subsection{Filières en pharmacie}

\paragraph{Répartition des primo-inscrits.}
Parmi les \num{3044}~primo-inscrits de 2025, 71,0~\% entrent en section~D
(adjoints d'officine et autres)~\cite{cnop2025}. La part de la section~D
était de 66~\% en 2015~\cite{cnop2016}.
\begin{itemize}
\item Les sections A, B, C et D sont privées par construction. Elles
  représentent 81,2~\%\derive{} du flux.
\item L'ambiguïté public / privé tient donc à G et H, soit 15,2~\%\derive{}
  du flux.
\end{itemize}

\paragraph{Parts privées de G et H.}
\begin{itemize}
\item Biologie (G) : 55,8~\%\derive{} de privé.
\item Établissements (H) : 28,6~\%\derive{} de privé selon l'Ordre, en
  inscriptions, avec des inscriptions multiples. La DREES donne 17~\%
  (axe~4), en personnes.
\end{itemize}
Le défaut proposé est la valeur de la DREES, cohérente avec le stock du
RPPS ; l'Ordre fournit la borne haute.

\paragraph{Internat.}
L'internat ouvre \num{746}~postes à la rentrée 2025~\cite{cnop2024}, soit
environ 14~\%\derive{} d'une promotion.

\paragraph{Diplômés non inscrits.}
Le taux de diplômés non inscrits à l'Ordre à trois ans était de 32,8~\% en
2015~\cite{cnop2016}. Aucune valeur récente n'a été vérifiée.

\subsection{Spécialisation en odontologie}

\paragraph{Postes de DES.}
Le concours national offre \num{121}~postes de DES en 2024--2025 : \num{52}
en orthopédie dento-faciale, \num{25} en chirurgie orale et \num{44} en
médecine bucco-dentaire~\cite{arrete2024odonto}. L'année précédente en
offrait \num{120}~\cite{arrete2023odonto}, et 2018 en offrait
\num{110}~\cite{cng2018}.

Rapportés aux \num{1322}~places en deuxième année quatre ans plus tôt, ces
postes représentent 9,2~\%\derive{} d'une promotion. Le décalage de quatre
ans est \emph{supposé}\suppose{}.

\paragraph{Stock de spécialistes.}
Le stock d'orthodontistes est stable : \num{2441} en 2019 et \num{2465} en
2022~\cite{oncd2022}. Les spécialistes, exemptés de régulation, forment donc
un sous-stock de 6~\%\derive{} qu'il faut séparer des omnipraticiens dans le
module de régulation.

\subsection{Abandon en études infirmières}

\paragraph{Le chiffre de 30~\%.}
Le chiffre de \enquote{30~\% d'abandon} vient d'une réponse ministérielle à
une question orale~\cite{qo101}. Cette réponse ne cite aucune source et
reconnaît l'absence de données.

\paragraph{Ce que mesure la DREES.}
\begin{itemize}
\item Les interruptions en première année atteignent 12,0 à 13,5~\% sur
  2022--2025~\cite{drees491}, toutes interruptions comprises.
\item Le taux d'interruption par promotion, toutes interruptions comprises,
  passe de 10,4~\% (2014) à 25,0~\% (2025)~\cite{drees491}. Deux ruptures de
  définition interviennent, en 2022 et en 2025.
\end{itemize}

\paragraph{Rendement de formation.}
Le rendement $\rho$ du tableau~\ref{tab:rho} est la quantité que le module
formation peut lire chaque année. Il est passé d'environ \num{0.84} à
\num{0.72}. La rupture coïncide avec l'entrée dans Parcoursup (2019) et
avec la hausse des quotas, qui gonfle le dénominateur avant les diplômes.

Notre calcul pour l'entrée 2018 ($\num{0.814}$\mesure) retrouve les 81~\%
publiés par la DREES~\cite{er1266}.

\begin{table}[htbp]
\centering\small
\caption{Rendement de la formation infirmière, $\rho = D(t+3)/I_1(t)$. Calcul\mesure{} sur la série longue de la DREES~\cite{drees491}, champ France métropolitaine et DROM.}\label{tab:rho}
\begin{tabular}{rrrrr}
\toprule
Entrée $t$ & $I_1(t)$ & $D(t+3)$ & $\rho$ & $1-\rho$ \\
\midrule
2009 & \num{31302} & \num{26447} & \num{0.845} & 15,5~\% \\
2012 & \num{30735} & \num{25888} & \num{0.842} & 15,8~\% \\
2015 & \num{31322} & \num{25757} & \num{0.822} & 17,8~\% \\
2018 & \num{30182} & \num{24557} & \num{0.814} & 18,6~\% \\
2019 & \num{32114} & \num{24379} & \num{0.759} & 24,1~\% \\
2020 & \num{32450} & \num{24215} & \num{0.746} & 25,4~\% \\
2021 & \num{35355} & \num{25249} & \num{0.714} & 28,6~\% \\
2022 & \num{36218} & \num{26020} & \num{0.718} & 28,2~\% \\
\bottomrule
\end{tabular}
\end{table}

\subsection{Sorties et entrée en exercice}

\paragraph{Ce qui n'est pas publié.}
La DREES ne publie pas les probabilités de cessation par âge~\cite{dd76}.
Elles se mesurent cependant sur le stock du RPPS par âge, jusqu'à la classe
ouverte des 65~ans et plus.

\paragraph{Valeurs publiées.}
\begin{itemize}
\item \textbf{Âge de retraite des kinésithérapeutes} : 66,8 à 67,3~ans en
  moyenne~\cite{cnomk2024}.
\item \textbf{Cumul emploi-retraite} : de 0,8~\%\derive{} des cotisants
  libéraux (sages-femmes) à 4,7~\%\derive{} (dentistes), données
  2015~\cite{cor2016}.
\item \textbf{Âge d'entrée en exercice} : 24~ans pour les sages-femmes,
  26~ans pour les dentistes, 27 à 29~ans pour les pharmaciens~\cite{dd76}.
\end{itemize}

\subsection{Un panel de politiques jouables}

\paragraph{Pourquoi ce panel.}
La finalité du projet est de laisser le décideur jouer, dans l'interface,
sur un large panel de politiques applicables. Chaque politique doit donc se
traduire en réglage d'un hyperparamètre du catalogue.

Le tableau~\ref{tab:leviers} couvre 14~politiques, regroupées en 12~lignes (la régulation des quatre professions tient en une ligne). Il donne pour chacune :
\begin{itemize}
\item le paramètre déplacé ;
\item la plage jouable proposée ;
\item la preuve disponible de l'effet ;
\item l'existence ou non d'un point d'entrée dans le code du jumeau à la
  date du prétirage.
\end{itemize}

\paragraph{Ce qu'il montre.}
\begin{itemize}
\item Huit politiques ont un point d'entrée dans le code.
\item Trois sont jouables dès maintenant : abandon en IFSI, places en
  formation, cumul emploi-retraite. Deux d'entre elles ont un défaut à
  corriger.
\item Les leviers retraite et régulation existent, mais sont suspendus
  après la seconde revue du jumeau.
\item Six politiques demandent un module encore absent (installation
  communale, filières) ou un flux séparé.
\end{itemize}
Le goulot du panel est donc le code, pas la documentation des paramètres.

\begin{table}[htbp]
\centering\scriptsize
\caption{Panel de politiques publiques jouables. Code : \textbf{oui} = point d'entrée actif ; \textbf{susp.} = module présent, levier suspendu ; \textbf{non} = module absent.}\label{tab:leviers}
\begin{tabular}{p{3.6cm}p{2.9cm}p{2.6cm}p{3.2cm}p{0.8cm}}
\toprule
Politique & Hyperparamètre & Plage jouable & Preuve de l'effet & Code \\
\midrule
Réguler l'installation (IDE, SF, MK, CD) par liste de communes & règle $E\le X$, $\theta_p$, $\psi_{p,d}$ & 0 à 100~\% des communes ; $\theta\in[\num{0.5};1]$\suppose & IDE seulement~\cite{irdes270} & susp. \\
Étendre la régulation aux dentistes spécialistes & exemption CD & oui / non & aucune & non \\
Règle des deux tiers en périphérie (IDE) & règle d'activité & oui / non & aucune & non \\
Obligation d'exercice en zone sous-dotée, étendue au-delà des MK & durée par profession & 0 à 3~ans\suppose & observable après 2027 & non \\
Réduire l'abandon en IFSI & $\rho_{\mathrm{IDE}}$, soit \num{0.056} à \num{0.105} par an\derive & niveaux 2009--2018 à 2022 & équivalent à +14~\% de places~\cite{er1319} & oui \\
Places en formation & capacité & $\times[\num{0.8};\num{1.2}]$ & $\pm$12 à 13~\% à 2050~\cite{dd76} & oui \\
Postes d'internat en pharmacie & $\varphi_G$, $\varphi_H$ & $\pm$20~\%\suppose & mécanique & non \\
Postes de DES en odontologie & part de spécialistes & 100 à 160\suppose & mécanique & non \\
Retenir les diplômés en pharmacie & taux d'inscription & 20 à 35~\%~\cite{cnop2016} & aucune & non \\
Plafonner ou encourager le cumul emploi-retraite & part en cumul & $\times[\num{0.5};2]$ (axe~5) & aucune propre & oui \\
Reporter l'âge de départ & décalage de retraite & 0 à 2~ans & +1~\% d'effectif par an~\cite{er1062,er1075} & susp. \\
Diplômés étrangers & flux d'entrée & $\times[\num{0.5};2]$\suppose & variante DREES~\cite{dd76} & partiel \\
\bottomrule
\end{tabular}
\end{table}

\subsection{Priorisation par l'analyse de sensibilité (axe~5)}

Le classement de Morris puis de Sobol de l'axe~5 n'est pas encore publié à
la date de ce prétirage ; il sera intégré à la version suivante. Son plan
déclaré permet déjà trois constats.
\begin{itemize}
\item Il inclut deux paramètres de ce catalogue : la part en cumul
  emploi-retraite, tirée entre $0,5$ et $2$ fois la valeur de la fiche, et
  le décalage de la retraite, de $0$ à $2$~ans.
\item Il exclut par construction la formation et la régulation. Le jumeau
  tourne au département, alors que le zonage est communal, et le vivier en
  formation n'alimente pas l'effectif.
\item Il ne pourra donc pas classer $\rho_{\mathrm{IDE}}$, $\theta_p$ ni
  $\psi_{p,d}$. Les sourcer plus finement reste à justifier par une analyse
  au niveau communal.
\end{itemize}

\section{Validation et menaces à la validité}

\paragraph{Ce qui est validé.}
Deux recoupements sont faits.
\begin{itemize}
\item Le rendement $\rho$ reproduit la valeur publiée par la DREES pour
  l'entrée 2018 (81,4~\%\mesure{} contre 81~\%~\cite{er1266}).
\item Le zonage data.gouv concorde avec l'extraction CartoSanté de l'axe~4.
\end{itemize}
L'effet de la régulation infirmière sert de cible au protocole N3 de
l'axe~3, qui n'est pas validé à ce jour.

\paragraph{Menaces.}
\begin{enumerate}
\item \textbf{Section et employeur.} Les filières $\varphi$ décrivent la
  section d'inscription à l'Ordre, pas l'employeur. Une partie de
  l'industrie ne s'inscrit pas, donc $\varphi_B$ est sous-estimé.
\item \textbf{Pas de suivi individuel en IFSI.} La DREES raisonne sur des
  promotions théoriques. La définition des interruptions change en 2022,
  puis le questionnaire en 2025.
\item \textbf{Textes non lus.}
  \begin{itemize}
  \item Les arrêtés de méthodologie du zonage (2020, 2024), cités par les
    ARS.
  \item Le texte au Journal officiel de l'avenant~7 des sages-femmes.
  \item L'avenant~11 des infirmiers (2026), qui pourrait modifier la règle.
  \end{itemize}
\item \textbf{Validité externe de l'effet IRDES.} Il porte sur le zonage de
  2011 et sur 2006--2016 ; rien de chiffré ne couvre l'après-2019.
\item \textbf{Millésimes anciens.} Le cumul emploi-retraite et le taux
  d'inscription des pharmaciens datent de 2015.
\end{enumerate}

\section{Reproductibilité}

\paragraph{Où sont les fiches.}
\begin{itemize}
\item Chaque fiche, avec ses liens et ses millésimes :
  \texttt{v2/docs/lab/fiches/param-*.md}.
\item La synthèse :
  \texttt{v2/docs/lab/reports/2026-11-14-axe1-hyperparametres.md}.
\end{itemize}

\paragraph{Ce qui se recalcule.}
\begin{itemize}
\item \textbf{Rendement $\rho$.} C'est le seul calcul propre. Il se
  reproduit avec le script de l'annexe~\ref{app:script}, à partir du
  fichier \enquote{Séries longues -- 1966-2025} de la DREES~\cite{drees491}.
\item \textbf{Parts de zonage.} Elles se recalculent sur
  \texttt{t-zonages-com.csv}~\cite{zonagesdatagouv}. Il faut sommer la
  population municipale par niveau des colonnes \texttt{zide},
  \texttt{zmk}, \texttt{zsf} et \texttt{zdent}.
\end{itemize}
Aucune graine n'intervient : aucun calcul n'est aléatoire.

\section{Conclusion}

\paragraph{Ce que l'axe~1 apporte.}
Il fournit au jumeau 21~paramètres dotés d'une valeur sourcée, mesurée ou
supposée explicite. Il désigne aussi six paramètres que personne ne mesure
encore.

\paragraph{Ce que cela change pour les autres axes.}
\begin{itemize}
\item \textbf{Axe~0, module de régulation.} Il doit intégrer :
  \begin{itemize}
  \item l'exemption des spécialistes dentaires ;
  \item la règle des deux tiers pour les infirmiers ;
  \item l'obligation d'exercice des kinésithérapeutes à partir de 2027 ;
  \item $\theta_p$ et $\psi_{p,d}$ comme paramètres de scénario.
  \end{itemize}
\item \textbf{Axe~0, module de formation.} Il lit $\rho_{\mathrm{IDE}}$ au
  lieu d'un abandon supposé.
\item \textbf{Axe~3.} Il gagne une seconde cible, l'IRDES, et deux
  expériences naturelles à construire : sages-femmes et kinésithérapeutes.
\end{itemize}

\paragraph{Questions ouvertes, avec leur coût.}
\begin{itemize}
\item Différence de différences pour les sages-femmes et les
  kinésithérapeutes : environ 3~h chacune.
\item Destination des entrants bloqués $\psi$, à partir des installations
  dans les communes limitrophes : environ 4~h.
\item Cessation par âge sur le RPPS : environ 1~h.
\item Lecture des arrêtés de méthodologie du zonage et de l'avenant~11 des
  infirmiers : environ 50~min.
\end{itemize}

\bibliographystyle{plainnat}
\bibliography{refs}
\appendix
\section{Script du rendement de formation infirmière}\label{app:script}
\begin{verbatim}
# rho_IDE = diplomes(t+3) / inscrits en 1re annee(t)
# Fichier : DREES jeu 491, "Series longues - 1966-2025.xlsx"
# Usage : python3 -I ide_rendement.py <series_longues.xlsx>
import sys
import pandas as pd

def serie(f, onglet):
    d = pd.read_excel(f, onglet, header=None)
    annees = [int(float(str(a).rstrip('*')))
              if str(a).rstrip('*').replace('.', '').isdigit() else None
              for a in d.iloc[4]]
    lib = "Infirmiers diplômés d'État"
    ligne = d[d[1].astype(str).str.strip() == lib].iloc[0]
    return {a: int(v) for a, v in zip(annees, ligne)
            if a and str(v).strip() not in ('-', 'nan')}

f = sys.argv[1]
dip, ins = serie(f, 'Diplomes_an'), serie(f, 'Inscrits1ere_an')
print('entree_t,inscrits_1A,diplomes_t+3,rho,1-rho')
for t in sorted(ins):
    if t + 3 in dip:
        r = dip[t + 3] / ins[t]
        print(f'{t},{ins[t]},{dip[t+3]},{r:.3f},{1-r:.1%}')
\end{verbatim}
\end{document}
